\chapter{Healthy and Active Children}

\section*{Definition and Overview}
Healthy and active children are those who eat well, move regularly, sleep enough, and develop physically, mentally, and socially in a positive way. Australian guidelines recommend that children and young people aged 5--17 accumulate at least 60 minutes of moderate to vigorous physical activity every day, and limit recreational screen time to no more than two hours a day. Healthy habits in childhood establish the foundation for lifelong health, preventing obesity, type 2 diabetes, heart disease, and poor mental health in later life. Parents and caregivers play the central role in modelling and supporting healthy eating, activity, sleep, and screen habits.

\section*{Epidemiology}
The health of Australian children has improved in many ways, but there are significant concerns. Around 1 in 4 Australian children and adolescents is overweight or obese, and physical activity levels are below the guidelines for many. Around 6 in 10 children do not meet the recommended 60 minutes of daily activity, and screen time is a growing concern, with many children exceeding recommended limits. The prevalence of childhood obesity has risen over recent decades, and it tracks into adulthood, making early healthy habits critical.

\section*{Aetiology and Pathophysiology}
Child health is shaped by nutrition, activity, sleep, and the family environment.

\begin{itemize}
    \item \textbf{Nutrition:} a balanced diet supports growth, development, immunity, and learning
    \item \textbf{Physical activity:} strengthens bones and muscles, supports heart health, and improves mood, sleep, and concentration
    \item \textbf{Sleep:} adequate sleep is essential for growth, learning, and emotional regulation
    \item \textbf{Screen time:} excessive screen time is linked to inactivity, poor sleep, and behavioural problems
    \item \textbf{Family environment:} children model the habits of their parents and caregivers
    \item \textbf{Growth and development:} regular checks track healthy growth and development
    \item \textbf{Prevention:} healthy habits prevent obesity, diabetes, and heart disease
\end{itemize}

\section*{Clinical Features}
Signs that a child is healthy and active include:

\begin{itemize}
    \item Steady growth along expected growth curves
    \item Age-appropriate development in motor, language, and social skills
    \item Daily physical activity, including active play, sport, and outdoor time
    \item A balanced diet with vegetables, fruit, whole grains, and protein
    \item Adequate sleep for the child's age
    \item Limited screen time, balanced with activity and sleep
    \item Good energy, concentration, and mood
    \item A healthy weight for the child's age and height
\end{itemize}

\section*{Differential Diagnosis}
Concerns in child health can have multiple causes.

\begin{itemize}
    \item Poor growth or weight gain, which may reflect diet, medical conditions, or feeding difficulties
    \item Overweight and obesity, which are influenced by diet, activity, and family factors
    \item Poor sleep, which can reflect habits, screen use, or medical conditions
    \item Behavioural and concentration problems, which may relate to sleep, diet, or development
    \item Delays in development, which require assessment
    \item Recurrent infections, which are common in childhood and usually normal
    \item Conditions affecting growth, such as hormonal or chronic disease
\end{itemize}

\section*{When to Seek Medical Care}
Children should attend routine health checks, including the health checks funded by Medicare at key ages, and immunisation visits. See a doctor for concerns about growth, development, weight, sleep, or behaviour, or for persistent symptoms.

\textbf{Call 000 (or local emergency services) for:}
\begin{itemize}
    \item A child with difficulty breathing, a seizure, or loss of consciousness
    \item A child who is difficult to wake or very unwell
    \item Signs of severe dehydration
    \item Any serious injury
    \item A child with fever and a non-blanching rash
\end{itemize}

\section*{Diagnosis and Investigations}
Health checks for children track growth and development.

\begin{itemize}
    \item \textbf{Growth measurement:} height, weight, and head circumference plotted on growth charts
    \item \textbf{Development checks:} milestones in movement, language, and social skills
    \item \textbf{Physical examination:} general health, vision, and hearing checks
    \item \textbf{Lifestyle review:} diet, activity, sleep, and screen time
    \item \textbf{Dental review:} from around age two
    \item \textbf{Immunisation check:} ensuring vaccinations are up to date
    \item \textbf{Referral:} to paediatric, dietetic, or other services when concerns arise
\end{itemize}

\section*{Management}
Supporting healthy and active children is a family and community effort.

\begin{itemize}
    \item \textbf{Modelling:} parents who eat well and are active have children who do the same
    \item \textbf{Family meals:} regular, shared meals support healthy eating
    \item \textbf{Active play:} encourage outdoor play, sport, and family activities
    \item \textbf{Limit screens:} set boundaries on screen time and keep screens out of bedrooms
    \item \textbf{Sleep routines:} consistent bedtimes and age-appropriate sleep
    \item \textbf{Healthy food environment:} keep healthy foods available and limit sugary drinks and treats
    \item \textbf{School and community:} active travel, school sport, and community programs
    \item \textbf{Professional support:} dietitians, paediatricians, and psychologists for specific concerns
\end{itemize}

\section*{Complications}
\begin{itemize}
    \item Obesity and its complications, including diabetes, joint problems, and low self-esteem
    \item Poor fitness and cardiorespiratory health
    \item Impaired sleep and concentration
    \item Behavioural and emotional problems
    \item Dental decay from sugary foods and drinks
    \item Missed developmental concerns without routine checks
    \item Long-term disease risk from unhealthy childhood habits
\end{itemize}

\section*{Prognosis and Outcomes}
Children who adopt healthy habits are more likely to maintain them into adulthood, with lower risks of obesity, diabetes, heart disease, and mental health problems. Physical activity in childhood improves fitness, bone health, academic performance, and wellbeing. Most concerns, including weight and sleep issues, respond well to early, family-centred intervention. The childhood years are the most powerful opportunity to establish lifelong health.

\section*{Prevention and Self-Care}
\begin{itemize}
    \item Model healthy eating and activity for your children
    \item Ensure at least 60 minutes of physical activity daily for children 5--17
    \item Limit recreational screen time to under two hours a day
    \item Provide a balanced diet with plenty of vegetables, fruit, and whole grains
    \item Limit sugary drinks and highly processed foods
    \item Establish consistent sleep routines
    \item Attend routine health checks and immunisations
    \item Encourage active play, sport, and outdoor time
    \item Make healthy choices the easy choices at home
    \item Seek help early for concerns about weight, development, or behaviour
\end{itemize}

\section*{Sources and Further Reading}
The following reputable sources provide further information for healthcare students and patients seeking reliable, current advice about healthy and active children.

\begin{itemize}
    \item Australian Government Department of Health and Aged Care. \textit{Physical activity guidelines for children.} https://www.health.gov.au/
    \item Healthdirect Australia. \textit{Healthy habits for children.} https://www.healthdirect.gov.au/
    \item Raising Children Network (Australian Government). \textit{Healthy eating and activity.} https://raisingchildren.net.au/
    \item Australian Institute of Health and Welfare. \textit{Child health data.} https://www.aihw.gov.au/
    \item World Health Organization. \textit{Physical activity in children.} https://www.who.int/
    \item NHS. \textit{Physical activity guidelines for children.} https://www.nhs.uk/
\end{itemize}
